% Construcción dodecafásica del funcional de retorno.
\section{Dodecaphase fundamental chain and hexad--octad functional}

The active incidence is constructed before any gravitational realization.
Let
\[
 H_{\mathrm{act}}=\{1,2,4,5,7,8\},
 \qquad
 \mathcal P_2(H_{\mathrm{act}})
 =\{A\subset H_{\mathrm{act}}:|A|=2\}.
\]
Then
\[
 |H_{\mathrm{act}}|=6,
 \qquad
 |\mathcal P_2(H_{\mathrm{act}})|=15,
\]
and the incidences before and after enlargement have cardinalities
\[
 \mathcal F_{90}=H_{\mathrm{act}}\times
 \mathcal P_2(H_{\mathrm{act}}),
 \qquad |\mathcal F_{90}|=90,
\]
\[
 \mathcal F_{120}=\{1,\ldots,8\}\times
 \mathcal P_2(H_{\mathrm{act}}),
 \qquad |\mathcal F_{120}|=120.
\]
These two cardinalities are outputs of the phase rule; they are neither angles nor
parameters supplied by the physical realization.

Denote by \(C_{12}=\mathbb Z/12\mathbb Z\) the temporal record of TPK and
by \(\partial_{\rm ret}\) its unique return seam per turn. In the
free module of oriented events, consider the fundamental chain
\begin{equation}
 c_{12}^{+}
 =\sum_{j\in C_{12}}[j,\mathcal F_{90}]
  -[\partial_{\rm ret},\mathcal F_{120}].
 \label{intsuccpass:eq:c12-fundamental-chain}
\end{equation}
The sum contains exactly the twelve sectors of the wheel; the second
term is unique because one turn has a single closure seam. Its sign
is the boundary sign. The bidirectional involution sends
\(c_{12}^{+}\) to \(c_{12}^{-}=-c_{12}^{+}\): it changes the orientation of the
chain without altering its absolute multiplicities or its incidence
types.

For \(m\in\{90,120\}\), let
\[
 S_m(q)=\frac{q^m}{1-q^{3m}}.
\]
The return readout \(\mathscr R\) is defined on the generators of
\eqref{intsuccpass:eq:c12-fundamental-chain} by
\[
 \mathscr R([j,\mathcal F_{90}])=S_{90},
 \qquad
 \mathscr R([\partial_{\rm ret},\mathcal F_{120}])=S_{120}.
\]
The first equality is constant on \(C_{12}\) by covariance under
phase rotation; the second has a different domain and acts only at the
enlarged seam. By linearity,
\begin{equation}
 \boxed{\mathscr R(c_{12}^{+})=12S_{90}-S_{120}.}
 \label{intsuccpass:eq:dodecaphase-forces-12minus1}
\end{equation}
Thus the coefficients \(12:-1\) are the image of the fundamental chain: twelve
phase cells and one oriented boundary. They are not selected through the
conventional value of the Barbero--Immirzi parameter, through a decimal
comparison, or by prescribing in advance the pole coefficient to be calculated
next.

\begin{theorem}[Pole coefficient of the dodecaphase chain]
The functional \(F(q)=\mathscr R(c_{12}^{+})(q)
=12S_{90}(q)-S_{120}(q)\), determined by the twelve sectors
and their oriented seam, has coefficient \(1/24\) for
\((1-q)^{-1}\). Its complex residue at \(q=1\) is \(-1/24\).
The pole coefficient is an output of the already constructed functional.
\end{theorem}

\begin{proof}
For a functional with a simple pole at \(q=1\), write
\[
 p_1(F):=\lim_{q\uparrow1}(1-q)F(q).
\]
The identity
\[
 1-q^{3m}=(1-q)\sum_{j=0}^{3m-1}q^j
\]
implies
\[
 p_1(S_m)
 =\lim_{q\uparrow1}\frac{q^m}{\sum_{j=0}^{3m-1}q^j}
 =\frac1{3m}.
\]
By linearity,
\[
 p_1\bigl(\mathscr R(c_{12}^{+})\bigr)
 =\frac{12}{270}-\frac1{360}
 =\frac{48-3}{1080}=\frac1{24}.
\]
These equalities are taken in \(\mathbb Q\). The complex residue
is the coefficient of \((q-1)^{-1}\), so that
\[
 \operatorname*{Res}_{q=1}\mathscr R(c_{12}^{+})(q)
 =-p_1\bigl(\mathscr R(c_{12}^{+})\bigr)=-\frac1{24}.
\]
\end{proof}

The Diophantine equation
\[
 \frac a{270}-\frac b{360}=\frac1{24}
 \quad\Longleftrightarrow\quad
 4a-3b=45
\]
remains a subsequent arithmetic certificate. On its own, it admits the
family \((a,b)=(12+3t,1+4t)\); the chain
\eqref{intsuccpass:eq:c12-fundamental-chain} eliminates that ambiguity before the calculation,
since changing \((12,1)\) alters the multiplicity of the orbit or its seam
and therefore no longer represents a fundamental TPK turn.

With the conjugate channels \(q_\pm\) already generated by the angular branch, one
finally obtains
\[
 \mathscr R(c_{12}^{+})(q_-)-
 \mathscr R(c_{12}^{+})(q_+)
 =12\Delta S_{90}-\Delta S_{120}.
\]
This is the internal input of the hexadic--octadic functional. Its evaluation
as the Barbero--Immirzi parameter and its insertion into the area operator are
subsequent publications: they do not participate in the selection of
\(c_{12}^{+}\), \(90/120\), \(12:-1\) or \(1/24\).

\paragraph{Falsifiers.}
The construction is refuted if the phase census does not yield twelve sectors, if
one turn contains more than one seam, if the incidences do not have
cardinalities \(90\) and \(120\), if the involution does not reverse the entire chain,
or if the exact coefficient of \((1-q)^{-1}\) in its image differs from \(1/24\).
